%ENS, 2001

\setcounter{equation}{0}
\setcounter{Q}{0}

Dans cet exercice, nous allons chercher représenter les courbes de solidus et liquidus pour un équilibre solide-liquide à partir d'un diagramme d'enthalpie libre d'un mélange.

\paragraph*{Enthalpie libre molaire $G_m$ du mélange}

Un système contient une mole de matière répartie suivant $1-x$ \si{\mole} de \ce{A} et $x$ \si{\mole} de \ce{B}, \ce{A} et \ce{B} étant dans le même état physique à la température $T$ considérée

\Q Exprimer l'enthalpie libre molaire $G_{m,1}$ du mélange en fonction de $\mu_A^\star$, $\mu_B^\star$ et de la fraction molaire $x$ du composé \ce{B} dans le mélange, dans le cas où \ce{A} et \ce{B} sont non miscibles dans leur état physique commun.

\Q Exprimer l'enthalpie libre molaire $G_{m,2}$ du mélange en fonction de $\mu_A^\star$, $\mu_B^\star$ et de la fraction molaire $x$ du composé \ce{B} dans le mélange, dans le cas où \ce{A} et \ce{B} sont miscibles et forment une solution solide idéale.

\Q Tracer les graphes des deux fonctions $G_{m,1}(x)$ et $G_{m,2}(x)$ sur la même figure. Quelles sont les valeurs de la dérivée de la fonction $G_{m,2}(x)$ lorsque $x \to 0$ et $x \to 1$ ?

\Q Montrer que l'intersection avec l'axe $x=0$ (resp. $x=1$) de la tangente en un point d'abscisse $x = x_0$ de la courbe $G_{m,2}(x)$ correspond au potentiel chimique $\mu_A$ (resp. $\mu_B$).

Exprimer alors la valeur de $\mu_A$ en fonction de $G_{m,2}(x_0)$, de $x_0$ et de $\left( \frac{\partial G_{m,2}}{\partial x} \right) (x = x_0)$.

\paragraph*{Relation entre les diagrammes $G_m (x)$ et les diagrammes isobares}
On envisage des mélanges à l'état solide et à l'état liquide. Les températures de fusion des corps purs \ce{A} et \ce{B} sont notées $T_A$ et $T_B$.

On étudie un premier mélange binaire dont la miscibilité des constituants est totale dans les deux phases liquide et solide. Les diagrammes des enthalpies libres du mélange à l'état liquide $G_{m,L}(x)$ et du mélange à l'état solide $G_{m,s}(x)$ à différentes températures sont donnés sur les figures (1) à (5) de la Fig. \ref{fig:02-01}.

Ces schémas représentent l'enthalpie libre du mélange $\ce{(1 - x) A + x B}$ pour $0 < x < 1$ à différentes températures : $T1 = 200 ~ \si{\celsius}$ ; $T_A = 150 ~ \si{\celsius}$ ; $T_2 = 100 ~ \si{\celsius}$ ; $T_B = 70 ~ \si{\celsius}$ ; $T_3 = 50 ~ \si{\celsius}$.

\Q Quelles sont les phases les plus stables aux différentes températures correspondant aux
diagrammes (1), (2), (4) et (5)? Justifier votre réponse.

\Q On étudie maintenant le diagramme (3) correspondant à la température $T_0 = 100 ~ \si{\celsius}$.

\begin{enumerate}
\item Montrer que construire la tangente commune aux deux courbes $G_{m,L}(x)$ et $G_{m,S}(x)$ est équivalent à établir l'égalité des potentiels chimiques de chaque constituant dans les phases dont les compositions $x_1$ et $x_2$ sont données par les points de contact de la tangente commune ?

\item Quelles sont les phases en présence selon les valeurs de $x$ ? Justifier votre réponse.

\item Quelle est l'expression de $G_m(x_0, x_1, x_2)$, enthalpie libre d'un mélange hétérogène de fraction molaire initiale $x_0$ telle que $x_1 < x_o < x_2$ ?
\end{enumerate}

\Q En déduire l'allure du diagramme $T = f(x)$ pour ce binaire solide-liquide.

\begin{figure}
\centering
\includegraphics[width=0.8\textwidth]{figures/Fig01}
\captionof{figure}{Représentation de l'enthalpie libre de la phase liquide et de la phase solide à différentes températures. La pression est fixée à \SI{1}{\bar}.}
\label{fig:02-01}
\end{figure}